

\section{Introduction}
\label{sec:rh_intro}
\begin{itemize}
\item \textbf{Realized variance within the trading day.}
  \begin{itemize}
  \item Realized variance, the sum of squared high-frequency returns, made volatility observable and forecastable
    \citep{AndersenBollerslevDieboldLabys2003, BarndorffNielsenShephard2002}, and the heterogeneous autoregressive
    (HAR) model of \citet{Corsi2009}, a regression of realized variance on its own averages over several past
    horizons, has become the benchmark for forecasting it.
  \item Within the trading day volatility follows a pronounced diurnal pattern \citep{AndersenBollerslev1997}, and
    pooled models that fit one forecasting equation across the intraday intervals, after a diurnal adjustment, have
    carried the approach to that horizon \citep{EngleSokalska2012,ZhangZhangCucuringuQian2024}. Like the HAR model,
    they forecast volatility from its own past; we ask what predictors beyond its lagged realized variances add.
  \end{itemize}
\item \textbf{Research question.} We investigate whether exogenous predictors improve pooled forecasts of the
  realized variance of front-month E-mini S\&P~500 futures on every regular-hours half-hour bar of the trading day.
  \begin{itemize}
  \item We call exogenous every predictor outside the HAR components of the target and the calendar dummies. The
    \rhdmNSeriesPooled{} exogenous series fall into \bkNBuckets{} groups drawn from the futures' own higher-order
    realized moments and trading activity; the intraday trading of the index's constituents, through their realized
    moments, liquidity and order-flow imbalance; option-implied volatility; the calendar of FOMC announcements; and
    alternative data, namely sentiment and attention on StockTwits and the net option vega bought by
    non-market-makers (volatility demand).
  \item Alternative data of this kind have mostly been used to forecast daily realized volatility
    \citep{AudrinoSigristBallinari2020,CaporinPoli2017}.
  \end{itemize}
\item \textbf{Design.}
  \begin{itemize}
  \item Each model class is estimated as one pooled model on the \rhePoolNBars{} regular-hours bars of the day,
    $\mathcal I_{\mathrm{rh}}$ (bars ending 10:00 to 16:00 ET), as in the intraday volatility models of
    \citet{EngleSokalska2012} and \citet{ZhangZhangCucuringuQian2024}.
  \item Its predictors, the HAR components of the target and trailing averages of every exogenous series, are
    computed on the futures' full grid of \rhdmNIntervals{} half-hour intervals, overnight included, so that the
    forecasts use the information that accumulates outside regular hours \citep{StroudJohannes2014}.
  \item For bar $i \in \mathcal I_{\mathrm{rh}}$ of day $t$ the forecast is
    \[
      \hat y_{t,i} = f_{m,\mathcal X,\mathcal I_{\mathrm{rh}}}(\bm x_{t,i}), \qquad
      \mathcal X \in \{\mathrm{HAR}, \mathrm{Full}\},
    \]
    for model classes $m$ from two families, the penalized linear models (ridge regression, the lasso and the
    elastic net) and the tree ensembles (LightGBM, XGBoost and the random forest). The exogenous series are the exact
    difference between the two predictor sets.
  \item Every model is re-estimated at every bar on a rolling window of \rhorWindowDays{} trading days
    ($\rhorWindowBars$ bars), with hyperparameters fixed before the evaluation span and then frozen
    (Section~\ref{sec:rh_methods}).
  \end{itemize}
\item \textbf{Benchmark and evaluation.}
  \begin{itemize}
  \item The primary benchmark is ordinary least squares on the HAR benchmark set, pooled on the same bars (the OLS
    row of Table~\ref{tab:rh_oos}); on this set the penalized linear models choose penalties close to zero and
    reproduce its forecast.
  \item Every forecast is compared with this benchmark, and each class's full predictor set with the HAR benchmark
    set of the same class, the paired comparison that isolates what the exogenous series add.
  \item Accuracy is the average QLIKE loss over the \rhorNIntB{} regular-hours bars of the \rhorDays{} trading days
    from \rhorSpanStart{} to \rhorSpanEnd, with Diebold--Mariano tests on the daily average losses
    \citep{DieboldMariano1995}. QLIKE weighs an error relative to the forecast and, like squared error, ranks
    forecasts consistently despite noise in the realized-variance proxy \citep{Patton2011}.
  \end{itemize}
\item \textbf{Contributions.}
  \begin{enumerate}
  \item \textbf{The exogenous predictors improve the forecast of every bar, beyond richer HAR models.}
    \begin{itemize}
    \item On every regular-hours bar the full predictor set lowers the average QLIKE loss relative to the OLS
      benchmark by \rhorVsOlsRedLinRng{}\% in the penalized linear models and by
      \rhorVsOlsRedTreeRng{}\% in the tree ensembles, and relative to the HAR benchmark set of its own class
      in all \rhorAfPairs{} model classes, each reduction statistically significant at the 5\% level
      (Section~\ref{sec:rh_results_oos}).
    \item A richer volatility model built from the same one-minute futures returns does not deliver this gain
      (Section~\ref{sec:rh_fi_ordinary}). The full set beats the HAR benchmark augmented with the futures' own
      realized moments significantly in \rhabNClassesBeatAddRm{} of the \rhabNClassesAddRm{} classes (QLIKE changes
      of \rhabLgbmFullVsAddRmAllPct\% for LightGBM and \rhabEnetFullVsAddRmAllPct\% for the elastic net), whereas the
      semivariance, realized-quarticity and leverage extensions of the HAR model
      \citep{PattonSheppard2015,BollerslevPattonQuaedvlieg2016,CorsiReno2012}, estimated together by OLS, change the
      benchmark's loss by \rhabOlsSharqlAllPct\%.
    \item The result holds at a window of \rhrbWinLong{} trading days (\rhrbNFullLowerSigBarLong{} of
      \rhrbNModelsLong{} classes; Section~\ref{sec:rh_results_robust}).
    \end{itemize}
  \item \textbf{The gain is concentrated, not dense and weak.}
    \begin{itemize}
    \item Two of the \bkNBuckets{} predictor groups carry most of it (Section~\ref{sec:rh_results_ablation}). The
      FOMC calendar comes first: added alone to the HAR benchmark, it changes the QLIKE loss of LightGBM by
      \rhabLgbmAddFomcAllPct\% and that of XGBoost by \rhabXgbAddFomcAllPct\%, against \rhabLgbmFullAllPct\% and
      \rhabXgbFullAllPct\% with their full sets, and removing it from the full set raises their loss by
      \rhabLgbmLoFomcAllPct\% and \rhabXgbLoFomcAllPct\%. The futures' realized moments come second (removing them
      raises LightGBM's loss by \rhabLgbmLoRmAllPct\%).
    \item Along a ladder of nested predictor sets ranked before the evaluation span, the \ladderTopK{}
      highest-ranked series retain \ladderTreeRetainRng{}\% of the trees' full-set gain, and the penalized
      linear models do best with \ladderLinKstarRng{} series (Appendix~\ref{app:rh_fi_ladder},
      Figure~\ref{fig:rh_ladder}). On the dense-versus-sparse question, the exogenous signal sits in a few series
      rather than spreading thinly over many.
    \item Beyond the FOMC calendar and the futures' realized moments, the other \rhbyFLNSeries{} series change no
      penalized linear model's loss significantly and that of LightGBM and XGBoost by
      \rhabLgbmFullVsAddRmFomcAllPct{} and \rhabXgbFullVsAddRmFomcAllPct\% (Section~\ref{sec:rh_fi_rest}), through a
      few series from several groups, chiefly the constituents' absolute returns and spreads, the VIX and
      StockTwits sentiment.
    \item StockTwits and the net vega demand add little even on the bars on which their values are live: there
      StockTwits changes
      the trees' loss by \rhabLgbmAddSentLiveSentPct\% (LightGBM) and \rhabXgbAddSentLiveSentPct\% (XGBoost), and
      volatility demand raises that of the linear models (\rhabRidgeAddVdLiveVdPct\% for ridge). Their share of the
      trees' attribution, \rhfipSentMin--\rhfipSentMax\% for StockTwits over the span (Table~\ref{tab:rh_fip_shares}),
      rises to \rhfilStwLiveMin--\rhfilStwLiveMax\% in their live years (Figure~\ref{fig:rh_fi_live}), but use is not
      accuracy.
    \end{itemize}
  \item \textbf{The gain holds under QLIKE, not under squared-error loss.}
    \begin{itemize}
    \item On every bar the full set has a higher mean squared error than the HAR benchmark in
      \rhcaNMseHigherBarSpan{} of the \rhorAfPairs{} classes (significantly in \rhcaNMseHigherSigBarSpan).
    \item A handful of days in the COVID-19 crash account for it: on every bar of \rhcaYearsOos, the \rhcaTopDays{}
      days with the largest excess squared error, all between \rhcaTopFirstBarOos{} and \rhcaTopLastBarOos, carry
      \rhcaLinTopShareBarOosRng\% of the full set's excess squared error in the penalized linear models and
      \rhcaTreeTopShareBarOosRng\% in the tree ensembles. Without the days from \rhcaCrashStart{} to \rhcaCrashEnd{}
      the full set has the lower mean squared error in \rhcaNMseLowerBarOosExCrash{} of the classes,
      \rhcaNMseLowerSigBarOosExCrash{} of them significantly.
    \end{itemize}
  \end{enumerate}
\item \textbf{Application to 0DTE straddles} (Section~\ref{sec:rh_econ}).
  \begin{itemize}
  \item On each day with a same-day SPXW expiration we buy or sell one zero-days-to-expiration (0DTE) at-the-money
    straddle at 15:30 ET as the forecast of the closing bar's realized variance lies above or below the variance
    implied by the straddle's price; the straddle pays on the index's move over that half hour, not on its realized
    variance, so the forecast of the half hour's realized variance signals which side to take
    \citep{EngleHongKaneNoh1993,NohEngleKane1994}.
  \item Over \rheSzDays{} expiration days the full-set strategies beat the unconditional short straddle,
    significantly in \rhemFullVsShortSig{} of the \rhorAfPairs{} classes, but beat the strategy of the HAR benchmark
    significantly in \rhemFullVsOlsHarSig{} of them; their gain over the short straddle is earned on the
    \rhedLongPctMin--\rhedLongPctMax\% of days on which they buy and rests on each strategy's \rhedNTop{} best days,
    and none of them carries a significant market beta.
  \end{itemize}
\item \textbf{Organization.} Sections~\ref{sec:rh_related} to~\ref{sec:rh_methods} cover the literature, the data,
  the predictors and the methods; Section~\ref{sec:rh_results} reports accuracy, where the gain falls (on FOMC
  statement days and ordinary days) and which predictors carry it, and Section~\ref{sec:rh_econ} the application;
  Section~\ref{sec:rh_conclusion} concludes. The online appendix holds the implementation details, the robustness
  checks and the models estimated on the closing half hour alone (Appendix~\ref{app:rh_closing}).
\end{itemize}


